% ===========================================================================
\section{Sutcliffe's problem}\label{sec:arith}
% ===========================================================================
Source: \cite{MathMag}, which uses Theorems~2 and~3 of \cite{Sut}; we prove these too.

\begin{definition}[digit solutions]\label{def:sut-solution}
\lean{TJK.Sutcliffe.IsSolution, TJK.Sutcliffe.HasSolution}
For an integer $n\ge2$ and nonnegative integers $a_h,\dots,a_0$ write
$(a_h,\dots,a_1,a_0)_n=a_hn^h+\dots+a_1n+a_0$. An \emph{$(h+1)$-digit solution for $n$} is a
collection of integers $n,a_h,\dots,a_0,k$ with
\[
  n\ge2,\quad k>1,\quad a_h>0,\quad a_0>0,\quad 0\le a_i\le n-1\ (0\le i\le h),\qquad
  k\,(a_h,\dots,a_1,a_0)_n=(a_0,a_1,\dots,a_h)_n .
\]
We say \emph{there is an $(h+1)$-digit solution for $n$} if such $a_i,k$ exist. Thus a 2-digit
solution is $(a,c,k)$ with $k(an+c)=cn+a$, and a 3-digit solution is $(a,b,c,k)$ with
$k(an^2+bn+c)=cn^2+bn+a$, where $1\le a,c\le n-1$, $0\le b\le n-1$, $k\ge2$.
\end{definition}

\begin{lemma}[basic inequalities]\label{lem:sut-basic}
\lean{TJK.Sutcliffe.two_digit_basic, TJK.Sutcliffe.three_digit_basic}
\uses{def:sut-solution}
For a 2-digit solution $(a,c,k)$ for $n$, and for a 3-digit solution $(a,b,c,k)$ for $n$, we have
$k<n$ and $a<c$.
\end{lemma}
\begin{proof}
In the 3-digit case $an^2+bn+c\ge n^2$ because $a\ge1$, and $cn^2+bn+a\le (n-1)(n^2+n+1)=n^3-1<n^3$;
so $kn^2\le k(an^2+bn+c)<n^3$ and $k<n$. Since $k\ge2$,
$cn^2+bn+a\ge 2(an^2+bn+c)>an^2+bn+c$, i.e.\ $(c-a)(n^2-1)>0$, so $c>a$. The 2-digit case is the
same with $an+c\ge n$, $cn+a\le n^2-1$ and $(c-a)(n-1)>0$.
\end{proof}

\begin{proposition}[Sutcliffe, Theorem 2, first half]\label{prop:sut-two-composite}
\lean{TJK.Sutcliffe.not_prime_of_two_digit}
\uses{def:sut-solution, lem:sut-basic}
If there is a 2-digit solution for $n$, then $n+1$ is not prime.
\end{proposition}
\begin{proof}
Let $k(an+c)=cn+a$. Modulo $p=n+1$ we have $n\equiv-1$, so $k(c-a)\equiv a-c$, i.e.\
$p\mid(k+1)(c-a)$. By Lemma~\ref{lem:sut-basic}, $0<c-a<n<p$ and $3\le k+1\le n<p$. If $p$ were
prime it would divide one of the two factors, which is impossible.
\end{proof}

\begin{proposition}[explicit solutions]\label{prop:sut-explicit}
\lean{TJK.Sutcliffe.explicit_two_digit, TJK.Sutcliffe.explicit_three_digit}
\uses{def:sut-solution}
Let $n\ge2$ and suppose $n+1=de$ with integers $d\ge3$ and $e\ge2$. Put $a=e-1$, $c=n-e$, $k=d-1$.
Then $(a,c,k)$ is a 2-digit solution for $n$ and $(a,n-1,c,k)$ is a 3-digit solution for $n$.
\end{proposition}
\begin{proof}
Digits: $a=e-1\ge1$; $e\le(n+1)/3\le n-1$ gives $c=n-e\ge1$; clearly $a,c\le n-1$; and $k=d-1\ge2$.
Using $n+1=de$:
\begin{align*}
k(an+c)&=(d-1)\bigl((e-1)n+n-e\bigr)=(d-1)e(n-1),\\
cn+a&=(n-e)n+e-1=(n-1)(n+1-e)=(n-1)(de-e)=(d-1)e(n-1),\\
k\bigl(an^2+(n-1)n+c\bigr)&=(d-1)\bigl((e-1)n^2+n^2-n+n-e\bigr)=(d-1)e(n^2-1),\\
cn^2+(n-1)n+a&=(n-e)n^2+n^2-n+e-1=(n^2-1)(n+1-e)=(d-1)e(n^2-1).\qedhere
\end{align*}
\end{proof}

\begin{lemma}[factorization]\label{lem:sut-factor}
\lean{TJK.Sutcliffe.exists_factor}
If $n\ge4$ and $n+1$ is not prime, then $n+1=de$ with $d\ge3$ and $e\ge2$.
\end{lemma}
\begin{proof}
Write $n+1=uv$ with $2\le u\le v$. Then $v\ge3$, since $v=2$ would force $n+1=4$. Take $d=v$, $e=u$.
\end{proof}

\begin{theorem}[Sutcliffe, Theorem 2]\label{thm:sut-sutcliffe2}
\lean{TJK.Sutcliffe.two_digit_iff}
\uses{prop:sut-two-composite, prop:sut-explicit, lem:sut-factor}
Let $n\ge4$. There is a 2-digit solution for $n$ if and only if $n+1$ is not prime.
\end{theorem}
\begin{proof}
One direction is Proposition~\ref{prop:sut-two-composite}; the other is Lemma~\ref{lem:sut-factor} with
Proposition~\ref{prop:sut-explicit}.
\end{proof}

\begin{theorem}[Sutcliffe, Theorem 3]\label{thm:sut-sutcliffe3}
\lean{TJK.Sutcliffe.three_digit_of_two_digit}
\uses{lem:sut-basic, lem:sut-factor, prop:sut-explicit, prop:sut-two-composite}
If there is a 2-digit solution for $n$, then there is a 3-digit solution for $n$.
\end{theorem}
\begin{proof}
By Lemma~\ref{lem:sut-basic}, $1<k<n$, so $n\ge3$. For $n=3$ we would have $k=2$ and $2(3a+c)=3c+a$,
i.e.\ $c=5a\ge5>n-1$; impossible. So $n\ge4$, $n+1$ is not prime by
Proposition~\ref{prop:sut-two-composite}, and Lemma~\ref{lem:sut-factor} with
Proposition~\ref{prop:sut-explicit} gives a 3-digit solution. (In $\N$-valued arithmetic the
expressions $n(a-1)$, $a-b+c$ below are to be read in $\Z$.)
\end{proof}

\begin{lemma}[{\cite[Lemma 1]{MathMag}}]\label{lem:sut-lemma1}
\lean{TJK.Sutcliffe.kaczynski_lemma1}
There are no natural numbers $a,c,n$ with
$n\ge2$, $a<c\le n-1$, $a+c\ge n+1$, and $n(a-1)+c\mid n(c-1)+a$.
\end{lemma}
\begin{proof}
Suppose there are, and write
\begin{equation}\label{eq:sut1}
k\bigl(n(a-1)+c\bigr)=n(c-1)+a,\qquad k\ge1 .
\end{equation}
(Here $a\ge2$ because $a\ge n+1-c\ge2$, so both sides are positive and $k\ge1$.) Modulo $n-1$ we have
$n\equiv1$, so $k(a-n+c)\equiv c-n+a$, i.e.\ $n-1\mid(k-1)(a-n+c)$. Since $a+c\ge n+1$, $a-n+c>0$. If
$k=1$ then $n(a-1)+c=n(c-1)+a$, so $c-a=n(c-a)$ and $n=1$, contrary to $n\ge2$. Hence $k>1$,
$(k-1)(a-n+c)>0$, and so $(k-1)(a+c-n)\ge n-1$. Since $c-n\le-1$, $(k-1)(a-1)\ge(k-1)(a+c-n)\ge n-1$.
Therefore
\[
k\bigl(n(a-1)+c\bigr)=nk(a-1)+kc\ge n(k-1)(a-1)+kc\ge n(n-1)+kc\ge nc+kc>n(c-1)+a,
\]
using $c\le n-1$ and $kc>0\ge a-n$. This contradicts \eqref{eq:sut1}.
\end{proof}

\begin{lemma}[{\cite[Lemma 2]{MathMag}}]\label{lem:sut-lemma2}
\lean{TJK.Sutcliffe.kaczynski_lemma2}
\uses{lem:sut-basic}
If $n=2$ or $n=3$, there is no 3-digit solution for $n$.
\end{lemma}
\begin{proof}
By Lemma~\ref{lem:sut-basic}, $1<k<n$, which is impossible for $n=2$. For $n=3$ we get $k=2$ and
$2(9a+3b+c)=9c+3b+a$, i.e.\ $0=17a+3b-7c\ge17+3b-14>0$, a contradiction.
\end{proof}

\begin{theorem}[Kaczynski {\cite{MathMag}}; Sutcliffe's conjecture]\label{thm:sut-main}
\lean{TJK.Sutcliffe.two_digit_of_three_digit}
\uses{lem:sut-lemma1, lem:sut-lemma2, lem:sut-basic, thm:sut-sutcliffe2}
Let $n\ge2$. If there is a 3-digit solution for $n$, then there is a 2-digit solution for $n$.
\end{theorem}
\begin{proof}
By Lemma~\ref{lem:sut-lemma2} we may assume $n>3$, and by Theorem~\ref{thm:sut-sutcliffe2} we may
assume that $p=n+1$ is prime. Let $(a,b,c,k)$ be a 3-digit solution:
\begin{equation}\label{eq:sut3}
k(an^2+bn+c)=cn^2+bn+a .
\end{equation}
Modulo $p$, $n\equiv-1$, so $k(a-b+c)\equiv c-b+a$ and $p\mid(k-1)(a-b+c)$. If $p\nmid a-b+c$ then
$p\mid k-1$; since $k-1>0$ this gives $k-1\ge p$, impossible because $k<n$
(Lemma~\ref{lem:sut-basic}). Hence $p\mid a-b+c$. Now
$-p<-n<3-n\le a-b+c\le 2n-2<2p$, so $a-b+c=ep$ with $e\in\{0,1\}$. Substituting $b=a+c-ep$,
\[
an^2+bn+c=a(n^2+n)+c(n+1)-epn=p\bigl((a-e)n+c\bigr),\qquad cn^2+bn+a=p\bigl((c-e)n+a\bigr),
\]
so \eqref{eq:sut3} becomes $k\bigl((a-e)n+c\bigr)=(c-e)n+a$. If $e=1$ then
$k\bigl(n(a-1)+c\bigr)=n(c-1)+a$; moreover $0\le b=a+c-p$ gives $a+c\ge n+1$, and $a<c\le n-1$ by
Lemma~\ref{lem:sut-basic}. This is impossible by Lemma~\ref{lem:sut-lemma1}. Hence $e=0$ and
$k(an+c)=cn+a$ with $1\le a,c\le n-1$ and $k>1$: a 2-digit solution.
\end{proof}

\begin{corollary}[{\cite{MathMag}}]\label{cor:sut-main}
\lean{TJK.Sutcliffe.three_digit_iff}
\uses{thm:sut-main, thm:sut-sutcliffe2, thm:sut-sutcliffe3}
If $n>3$, there is a 3-digit solution for $n$ if and only if $n+1$ is not prime.
\end{corollary}
\begin{proof}
If $n+1$ is not prime, Theorem~\ref{thm:sut-sutcliffe2} gives a 2-digit solution and
Theorem~\ref{thm:sut-sutcliffe3} a 3-digit one. Conversely a 3-digit solution gives a 2-digit solution
(Theorem~\ref{thm:sut-main}), and then $n+1$ is not prime (Theorem~\ref{thm:sut-sutcliffe2}).
\end{proof}

\begin{remark}
\cite{MathMag} notes that a similar result was obtained independently by Prasert Na Nagara.
\end{remark}
